\documentclass[11pt,a4paper]{article}

% ---------------------------------------------------------
% PACKAGES
% ---------------------------------------------------------
\usepackage[utf8]{inputenc}
\usepackage[T1]{fontenc}
\usepackage{lmodern}
\usepackage{amsmath, amssymb, amsfonts, bm}
\usepackage{geometry}
\usepackage{graphicx}
\usepackage{hyperref}
\usepackage{authblk}
\usepackage{cite}
\usepackage{physics}
\usepackage{mathtools}
\usepackage{tensor}
\usepackage{enumitem}

\geometry{margin=2.4cm}

% ---------------------------------------------------------
% TITLE
% ---------------------------------------------------------
\title{\textbf{Companion III (v49.1)}\\[0.2cm]
\textbf{Relaxation Dynamics in a CPT-Symmetric Bipartite Universe}\\[0.2cm]
\large Foundations of Relaxation-Driven Cyclic Cosmology (RDCC)}

\author[1]{Michael Lehmann}
\affil[1]{\small Independent Researcher, Relaxation-Driven Cyclic Cosmology (RDCC)\\
\texttt{mi.lehmann@gmx.de}}

\date{September 2026}

\begin{document}
\maketitle

% ========================================================
% ABSTRACT
% ========================================================

\begin{abstract}
This Companion develops the relaxation sector of RDCC in its corrected and
fully EFT-coherent v49.1 form. The relaxation parameter $\alpha(t)$ arises
from the ultraviolet-suppressed inter-sector operator
$T^{(+)}_{\mu\nu}T^{(-)\mu\nu}$, which is identified in v49.1 as a
dimension-8 operator. This correction propagates through the infrared dynamics,
the RG flow, and the Lindblad-like effective dissipator. We derive the global
definition of $\alpha(t)$, establish its infrared fixpoint
$\alpha_{\mathrm{IR}}$, and show how relaxation controls all deviations from
$\Lambda$CDM. This Companion provides the mathematical and physical foundation
for the thermodynamic sector (Companion IV), the ultraviolet sector (Companion X),
and the tensor sector (Companion XIV).
\end{abstract}

\tableofcontents
\newpage
% ========================================================
% 1 INTRODUCTION
% ========================================================

\section{Introduction}

Relaxation-Driven Cyclic Cosmology (RDCC) describes the Universe as a globally
coherent CPT-symmetric quantum state defined on a bipartite Hilbert space,
\begin{equation}
\mathcal{H} = \mathcal{H}_{(+)} \otimes \mathcal{H}_{(-)}.
\end{equation}
The visible and mirror sectors arise as reduced density matrices of a single
global state. Their thermodynamic arrows of time point in opposite directions,
ensuring global CPT invariance.

The relaxation sector quantifies how the visible sector responds to the
ultraviolet-suppressed coupling between both sectors. This coupling is encoded
in the inter-sector operator
\begin{equation}
\mathcal{O}_{\mathrm{int}} = T^{(+)}_{\mu\nu} T^{(-)\mu\nu},
\end{equation}
which couples the energy-momentum tensors of both sectors through the global
gravitational field.

In earlier versions of RDCC (up to v48.0), this operator was implicitly treated
as a dimension-6 operator. In v49.1, we correct this and consistently treat
$\mathcal{O}_{\mathrm{int}}$ as a dimension-8 operator. This correction is
crucial for EFT coherence and propagates through the entire relaxation sector.

The goal of this Companion is to derive the corrected relaxation dynamics,
establish the global definition of $\alpha(t)$, determine its infrared fixpoint
$\alpha_{\mathrm{IR}}$, and present the Lindblad-like effective structure of
the dissipator.
% ========================================================
% 2 DIMENSION-8 INTER-SECTOR OPERATOR
% ========================================================

\section{Dimension-8 Inter-Sector Operator}

The energy-momentum tensor has mass dimension
\begin{equation}
[T_{\mu\nu}] = 4.
\end{equation}
Therefore, the inter-sector operator has mass dimension
\begin{equation}
[T^{(+)}_{\mu\nu} T^{(-)\mu\nu}] = 8.
\end{equation}
The corresponding term in the effective action is
\begin{equation}
S_{\mathrm{int}}
=
\int d^4x \sqrt{-g} \,
\frac{1}{M^4} \,
T^{(+)}_{\mu\nu} T^{(-)\mu\nu}.
\label{eq:dim8action}
\end{equation}

This correction ensures full EFT consistency. The scale $M$ parametrizes the
strength of the coupling and is assumed to be much larger than all cosmological
energy scales.
% --------------------------------------------------------
% 2.1 Global Definition of alpha(t)
% --------------------------------------------------------

\subsection[Global Definition of alpha(t)]{Global Definition of $\alpha(t)$}

The relaxation parameter $\alpha(t)$ is defined globally on the bipartite
Hilbert space. It is not a sectoral parameter; instead, it quantifies the
infrared response of the global state to the ultraviolet-suppressed coupling.

From Eq.~\eqref{eq:dim8action}, we obtain the scaling
\begin{equation}
\alpha(t) \propto \frac{1}{M^4}.
\end{equation}

The visible sector measures $\alpha(t)$ only indirectly through gravitational
observables. The mirror sector measures it through its thermodynamic arrow.
Both measurements correspond to projections of the same global parameter.
% ========================================================
% 3 RG FLOW AND INFRARED FIXPOINT
% ========================================================

\section{RG Flow and Infrared Fixpoint}

The effective coupling saturates in the infrared due to the RG flow,
\begin{equation}
\dot{\alpha}(t) = - \Gamma_{\mathrm{IR}} \, \alpha(t),
\end{equation}
where $\Gamma_{\mathrm{IR}}$ is the infrared relaxation rate. The solution is
\begin{equation}
\alpha(t)
=
\alpha_{\mathrm{IR}}
+
\left( \alpha_0 - \alpha_{\mathrm{IR}} \right)
e^{-\Gamma_{\mathrm{IR}} t}.
\end{equation}

The infrared fixpoint is defined as
\begin{equation}
\alpha_{\mathrm{IR}} = \lim_{t \to \infty} \alpha(t).
\end{equation}

This fixpoint controls all deviations from $\Lambda$CDM:
\begin{align}
n_s - 1 &\sim -2\alpha_{\mathrm{IR}}, \\
\Delta N_{\mathrm{eff}} &\sim \alpha_{\mathrm{IR}}, \\
\Delta(f\sigma_8) &\sim \alpha_{\mathrm{IR}}, \\
\Omega_{\mathrm{GW}} &\sim \alpha_{\mathrm{IR}}^2.
\end{align}
% ========================================================
% 4 LINDBLAD-LIKE EFFECTIVE DISSIPATOR
% ========================================================

\section{Lindblad-Like Effective Dissipator}

The ultraviolet-suppressed coupling induces an effective dissipator in the
visible sector. The global state evolves according to
\begin{equation}
\dot{\rho}_{(+)} =
- i [H_{(+)}, \rho_{(+)}]
+ \mathcal{D}_{\mathrm{eff}}[\rho_{(+)}],
\end{equation}
where the dissipator has the Lindblad-like form
\begin{equation}
\mathcal{D}_{\mathrm{eff}}[\rho]
=
\sum_{n} \gamma_n
\left(
L_n \rho L_n^\dagger
- \frac{1}{2} \{ L_n^\dagger L_n, \rho \}
\right).
\end{equation}

The coefficients $\gamma_n$ scale as
\begin{equation}
\gamma_n \propto \frac{1}{M^4},
\end{equation}
consistent with the dimension-8 suppression.

This structure explains the thermodynamic asymmetry between the sectors and
provides the microscopic origin of the relaxation dynamics.
% ========================================================
% 5 CONCLUSION
% ========================================================

\section{Conclusion}

This Companion develops the corrected and fully EFT-coherent relaxation sector
of RDCC. The identification of the inter-sector operator as a dimension-8
operator restores ultraviolet consistency and provides the microscopic origin
of the relaxation parameter $\alpha(t)$. The global definition of $\alpha(t)$,
its infrared fixpoint $\alpha_{\mathrm{IR}}$, and the Lindblad-like effective
structure of the dissipator form the backbone of RDCC's deviation from
$\Lambda$CDM.

These results propagate into the thermodynamic sector (Companion IV), the
ultraviolet sector (Companion X), and the tensor sector (Companion XIV). The
relaxation sector is therefore central to the physical interpretation and
observational predictions of RDCC.

\end{document}
